\section{Dataset and LOCO Protocol}
\label{sec:dataset}

\subsection{Sen1Floods11}

We use Sen1Floods11 \cite{bonafilia2020}, which provides $512\times512$
Sentinel-1 chips with labels from two sources. The hand-labeled subset
covers 11 flood events named by country or region, including the Mekong.
It contains 446 chips. Of these, 431 are in the standard split and 15
Bolivia chips are held out separately by design. The paired comparison
uses the ten standard groups and leaves out the separate Bolivia hand
holdout. The standard split has 252 training, 89 validation and 90 test
chips. The training share is 58.47\% of 431.

A much larger weakly labeled subset has labels derived automatically from
Sentinel-2 imagery. Bonafilia et al. report 4,385 weakly labeled chips.
The local manifests contain 4,384, with 3,946 train and 438 validation
entries across 12 filename groups, including Bolivia and Colombia. The
reason for the one-chip difference is unknown.

Water makes up 9.51\% of labeled pixels with finite SAR in the official
hand training split and 12.30\% in the official test split. The historical
10.6\% figure describes the full hand inventory. The labels mark water
present on the flood date, including both flood and persistent water.
This determines which model output can be scored (Sec.~\ref{sec:fusion}).

\subsection{Leave-One-Event-Out Protocol}
\label{sec:protocol}

For each planned fold, the reconstructed manifests remove the named
event from the hand and weak training and validation roles. Hand training
then uses 211--242 chips and hand validation uses 75--85. Three planned
seeds (42, 43 and 44) for the ten standard groups give 30 event--seed
combinations. The surviving runs cover only part of this design
(Sec.~\ref{sec:loco}).

The October reconstruction froze completed runs and their
checkpoints with the best validation IoU before paired inference. Committed
manifests, file hashes, run paths and validation histories support the
reconstruction of historical membership and checkpoint selection.
Runtime records of actual file reads and full initializer lineage were
unavailable, including the USA retry's exact pretraining initializer.

\subsection{Seed-Specific Paired Difference}
\label{sec:paired}

The original comparison scored the generalist across all 431 standard
hand chips, including the 252 used for hand training and 89 used for
validation. This mixed training, checkpoint-selection and test performance.
We now compare an event-excluded checkpoint for seed $s$ with a fixed
generalist of the same architecture and define
\begin{equation}
\Delta(e,s)=\operatorname{IoU}(\mathrm{LOCO}_{e,s},T_e)
             -\operatorname{IoU}(G_{42},T_e).
\label{eq:paired}
\end{equation}
Here $T_e$ contains only event $e$'s official test chips and $G_{42}$ is
the same seed-42 generalist checkpoint throughout. We use the same valid
mask in both arms, excluding ignored labels and non-finite SAR, and score
sigmoid probabilities strictly greater than 0.5 without test-time
augmentation or fusion. For each arm, water-class IoU is
$TP/(TP+FP+FN)$, using confusion counts pooled over the valid pixels in
$T_e$. A negative difference means the event-excluded checkpoint scored
lower on these inputs.

Identical test inputs control what is scored. Training budget,
initialization and validation composition still differ. The historical
Bolivia difference of $+0.141$ compared two models whose hand training
excluded Bolivia. It did not measure the effect of changing the
training budget. A later
size-matching pilot covered four events with one seed and one replacement
draw each. It provides descriptive results for the particular datasets
used. The accompanying reproduction package records the individual
runs and their design limits.
